\chapter*{Preface}
Volume~6 built the stage: real numbers, limits, measures and integrals. This book is about what happens on it when functions are
treated as points. A function becomes a vector in a space of infinite dimension, and the old words of linear algebra, basis,
dual, eigenvalue, projection, return with new meanings and new surprises. Infinite dimension is not just ``more of the same'':
the unit ball stops being compact, linear maps can be discontinuous, and an operator can have no eigenvalue at all.

The first two chapters are Banach and Hilbert spaces: Hahn--Banach and the three great
consequences of Baire's theorem, then Pythagoras for functions. Spectral Theory opens with matrices, the concrete case: duality and volume, eigenvalues, inner products, the spectral theorem for matrices, singular values and graphs. Cuts, expansion and the Riemann hypothesis for graphs come next, then the semicircle and the sine kernel, still finite spectra. Matrix inequalities and linear dynamics stay at that scale. Operators and the spectral theorem for compact self-adjoint operators follow, and one hears the shape of a vibrating string.

The next chapters are harmonic analysis, the art of writing a function as a superposition of waves. The Fourier transform
on $\R^d$, with Plancherel, Poisson summation, the uncertainty principle and Shannon's sampling theorem; distributions, which give
a derivative to every function, and Sobolev spaces, where differential equations find their solutions; Fourier analysis on
finite abelian groups, where it becomes the fast Fourier transform and proves quadratic reciprocity; and Haar measure on locally
compact groups and Pontryagin duality, with the $p$-adic numbers as a self-dual group. The adèles, which see $\R$ and every $\Q_p$ at once, are built in Volume~13.

The last two chapters return to the complex plane: Cauchy's theorem in its final form, using the winding number and the Jordan curve theorem of Volume~7, the growth of
entire functions and Mellin inversion. The Riemann mapping theorem is proved in Volume~9. Tauberian theorems then recover the behaviour of a
sequence from its generating function, and lead once more to the prime number theorem.

The book assumes Volume~6, the Fourier series and complex analysis of Volume~4, and the groups of Volume~2. Each chapter ends with
exercises, each with a hint and a complete solution.
